%=============================== Introducción =================================
% Estructura sugerida: contexto y motivación, planteo del problema, objetivos,
% contribuciones y organización del documento.

La \ac{IA} y, en particular, el \ac{PLN} han transformado la manera en que
se procesa la información textual. \todo{Contextualizar el área de la tesis.}

\section{Motivación y planteo del problema}
\todo{Describir el problema de investigación y por qué es relevante. Plantear
la brecha (gap) que la tesis viene a cubrir.}

\section{Objetivos}
\subsection{Objetivo general}
\todo{Enunciar el objetivo general en una o dos oraciones.}

\subsection{Objetivos específicos}
\begin{itemize}
  \item \todo{Objetivo específico 1.}
  \item \todo{Objetivo específico 2.}
  \item \todo{Objetivo específico 3.}
\end{itemize}

\section{Contribuciones}
Las principales contribuciones de esta tesis son:
\begin{enumerate}
  \item \todo{Contribución 1.}
  \item \todo{Contribución 2.}
\end{enumerate}

\section{Organización de la tesis}
El resto del documento se organiza como sigue. El \cref{cap:estado-del-arte}
revisa los antecedentes; el \cref{cap:marco-teorico} presenta el marco teórico;
el \cref{cap:metodologia} describe la metodología; los
\cref{cap:experimentos,cap:resultados} detallan los experimentos y sus
resultados; el \cref{cap:discusion} los discute; y el \cref{cap:conclusiones}
cierra con las conclusiones y el trabajo futuro.
